\documentclass[12pt,french,oneside]{report}
%%%%%%%%%%%%%%%%%%%%%%%%%%%%%%%%%%%%%%%%%%%%%%%%%%%%%%%%%%%%%%%%%%%%%%%%%%%%%%%
\input{preambule_2019}
%%%%%%%%%%%%%%%%%%%%%%%%%%%%%%%%%%%%%%%%%%%%%%%%%%%%%%%%%%%%%%%%%%%%%%%%%%%%%%%
\portrait
%\paysage
%%%%%%%%%%%%%%%%%%%%%%%%%%%%%%%%%%%%%%%%%%%%%%%%%%%%%%%%%%%%%%%%%%%%%%%%%%%%%%%
\renewcommand{\mp}[1]{\marginpar{\center \textbf{#1}}}
\newcommand{\classe}{MS 1}
\newcommand{\BTS}{MAINTENANCE DES SYSTÈMES}
%%%%%%%%%%%%%%%%%%%%%%%%%%%%%%%%%%%%%%%%%%%%%%%%%%%%%%%%%%%%%%%%%%%%%%%%%%%%%%%%%%%%%%
%entête bac blanc
\newcommand{\bacblanc}{
\thispagestyle{empty}

\includegraphics[scale=0.3]{logo_academie} \hfill 
\includegraphics[scale=0.2]{logo_augustin_thierry}
\normalsize

\vspace*{1em}

\hrule

\vspace*{1em}
\begin{center}
\LARGE
\begin{bf}
BTS \BTS\\
\vspace{1em}
\decofourleft~CCF de MATH\'EMATIQUES \decofourright\\
\vspace{1em}
Situation d'évaluation n\degres 1\\
\vspace{1em}
Sujet n\degres 1\\
\vspace{1em}
Vendredi 12 mai 2018\\
\vspace{1em}
Durée de l'épreuve : 55 min
\end{bf}
\end{center}
\normalsize

\vspace*{1em}


\hrulefill


\vspace*{1em}

\textbf{L'usage de la calculatrice ou des logiciels installés sur les ordinateurs est autorisé tout comme celui du formulaire établi ensemble.}

\textbf{Le sujet comporte deux exercices indépendants.}

\vspace*{1em}


\hrulefill


\vspace*{1em}

\textbf{COMP\'ETENCES EVALU\'EES}

\begin{center}
\small
\setstretch{1.2}
\begin{tabular}{m{0.5cm}>{\bf}m{4.5cm}m{10cm}}
\hline C1 & S'informer & Rechercher, extraire et organiser. l'information\tabularnewline
\hline C2 & Chercher & Proposer une méthode de résolution. Expérimenter, tester, conjecturer.\tabularnewline
\hline C3 & Modéliser & Représenter une situation ou des objets du monde réel. Traduire un problème en langage mathématique.\tabularnewline
\hline C4 & Raisonner, argumenter & Déduire, induire, justifier ou démontrer un résultat. Critiquer une démarche, un résultat.\tabularnewline
\hline C5 & Calculer, illustrer, mettre en \oe uvre une stratégie & Calculer, illustrer à la main ou à l'aide d'outils numériques, programmer.\tabularnewline
\hline C6 & Communiquer & Rendre compte d'une démarche, d'un résultat, à l'oral ou à l'écrit. Présenter un tableau, une figure, une représentation graphique.\tabularnewline \hline
\end{tabular}
\end{center}
\newpage
}
%%%%%%%%%%%%%%%%%%%%%%%%%%%%%%%%%%%%%%%%%%%%%%%%%%%%%%%%%%%%%%%%%%%%%%%%%%%%%%%
\tikzset{domaine/.style 2 args={domain=#1:#2}}
%%%%%%%%%%%%%%%%%%%%%%%
%% DEBUT DU DOCUMENT %%
%%%%%%%%%%%%%%%%%%%%%%%

\begin{document}
\selectlanguage{french}
\setlength\parindent{0mm}
\bacblanc 		%entête bac blanc

\renewcommand \footrulewidth{0.2pt}%
\renewcommand \headrulewidth{0.2pt}%
\pagestyle{fancy}
\fancyhead{Nom, Prénom :\dotfill Classe : STS MS 2
\vspace{6pt}}
\pieddepage{CCF BTS MS 2}{\thepage / \pageref{LastPage}}{Année 2018-2019}

%%%%%%%%%%%%%%%%%%%%%%%Page 1%%%%%%%%%%%%%%%%%%%%%%%%%%%%%%%%%%%%%%%%
\begin{spacing}{1.2}
%%%%%%%%%%%%%%%%%%%%%%%%%%%%%%%%%%%%%%%%%%%%%%%%%%%%%%%%%%%%
%%%%%%%%%%%%%%%%%%%%%%%%%%%%%%%%%%%%%%%%%%%%%%%%%%%%%%%%%%%%%%%%%%%%%%%%%
%%%%%%%%%%%%%%%%%%%%%%%% exercice 1 %%%%%%%%%%%%%%%%%%%%%%%%%%%%%%%%%%%%%
\label{exo1}
%%%%%%%%%%%%%%%%%%%%%%%%%%%%%%%%%%%%%%%%%%%%%%%%%%%%%%%%%%%%%%%%%%%%%%%%%
\begin{Exercice}[]%


\end{Exercice}

%%%%%%%%%%%%%%%%%%%%%%%%%%%%%%%%%%%%%%%%%%%%%%%%%%%%%%%%%%%%%%%%%%%%%%%%%
%%%%%%%%%%%%%%%%%%%%%%%% exercice 2 %%%%%%%%%%%%%%%%%%%%%%%%%%%%%%%%%%%%%
\label{exo2}
%%%%%%%%%%%%%%%%%%%%%%%%%%%%%%%%%%%%%%%%%%%%%%%%%%%%%%%%%%%%%%%%%%%%%%%%%
\newpage
%\medskip
\begin{Exercice}[( points)]%



\end{Exercice}




%%%%%%%%%%%%%%%%%%%%%%%%%%%%%%%%%%%%%%%%%%%%%%%%%%%%%%%%%%%%
%%%%%%%%%%%%%%%%%%%%%%%%%%%%%%%%%%%%%%%%%%%%%%%%%%%%%%%%%%%%
\end{spacing}

%%%%%%%%%%%%%%%%%%%%%%%%%%%%%%%%%%%%%%%%%%%%%%%%%%%%%%%%%%%%
%%%%%%%%%%%%%%%%%%%%%
%% FIN DU DOCUMENT %%
%%%%%%%%%%%%%%%%%%%%%
\end{document}